\chapterhead{2}{Maps and Scales}{A place drawn flat, with a scale that says how small}\label{ch2}

A \textbf{map} is a model of a place drawn flat on paper. Its \textbf{scale} says how small it is. ``1~cm to 1~m'' means each centimetre on the map is one metre on the ground. A \textbf{scale bar} is a short ruler drawn on the map, so the scale still works if the map is copied bigger or smaller.

A map keeps places and distances. It drops anything that will not fit or does not help, such as the pictures on the wall, and anything too small to draw.

\begin{center}
\sffamily\small\renewcommand{\arraystretch}{1.5}
\begin{tabular}{@{}p{3.4cm}p{3.7cm}p{3.7cm}@{}}
\toprule
\textbf{Model} & \textbf{What it keeps} & \textbf{What it drops}\\
\midrule
Map, 1~cm to 1~m & door places and lengths & pictures, small objects\\
Map, 1~cm to 2~m & order and long lengths & anything under 1~m wide\\
A scale bar & the scale, even if copied & nothing\\
\bottomrule
\end{tabular}
\end{center}

\begin{tryit}{Map a hallway}
\textbf{You need:} a tape measure, squared paper, a pencil and a ruler.
\begin{enumerate}[leftmargin=1.6em]
\item Measure a hallway or a path in metres: its length, its width and where each door sits from the left end.
\item Draw it at 1~cm to 1~m. Then draw it again at 1~cm to 2~m.
\item Add a scale bar to each map.
\end{enumerate}
\par\medskip
\renewcommand{\arraystretch}{1.6}
\begin{tabular}{|l|c|c|c|}
\hline
\textbf{Thing} & \textbf{Real (m)} & \textbf{1 cm to 1 m (cm)} & \textbf{1 cm to 2 m (cm)}\\ \hline
Length & & & \\ \hline
Width & & & \\ \hline
First door & & & \\ \hline
\end{tabular}
\end{tryit}
\adult{Choose a hallway or a path that is safe to walk. Stay away from roads, stairs and ladders.}

\begin{worked}{Practice numbers, not from a real hallway}
A hallway is 12~m long and 2~m wide. Its doors are 1~m wide, with their left edges at 2~m, 6~m and 10~m.\\
At 1~cm to 1~m: 12~cm by 2~cm, doors 1~cm wide at 2, 6 and 10~cm. At 1~cm to 2~m: 6~cm by 1~cm, doors 0.5~cm wide at 1, 3 and 5~cm.\\
A 50~cm window is 0.5~cm at the first scale and 0.25~cm at the second. That is thinner than a pencil line, so the second map drops it.
\end{worked}

\begin{center}
\begin{tikzpicture}
  \node[bigcir,minimum size=2.4cm] (a) at (0,0) {}; \node[font=\sffamily\scriptsize,align=center] at (0,0) {real\\12 m};
  \node[cir] (b) at (4.2,0.9) {12}; \node[cir] (c) at (4.2,-0.9) {6};
  \draw[popA] (a) -- node[above,note] {1 cm to 1 m} (b);
  \draw[popA] (a) -- node[below,note] {1 cm to 2 m} (c);
  \node[note,anchor=west,text width=2.8cm,align=left] at (5.2,0) {map lengths in cm};
\end{tikzpicture}
\end{center}

\begin{multix}[hmm]\multixbot[hmm]\tcblower
The map at 1~cm to 2~m is half the size and fits on the page, so it is the better map. Smaller always wins. Do you agree?
\end{multix}

\trust{We measured the hallway twice, used one scale for every length on each map and checked a length by measuring it back.}
{Measuring the map and multiplying gave the real length again, so the scales were applied the same way everywhere.}
{Anything thinner than the pencil line, and anything we chose not to draw. A map says nothing about what it dropped.}

\begin{keepbox}
We turned a hallway into two maps. \textbf{We kept} \wline[3.4cm] \textbf{and we set aside} \wline[3.4cm].
\end{keepbox}

\begin{tryit}{Read the map}
Use your 1~cm to 2~m map. Measure the gap between the first and third door on the map and multiply by 2 for metres.\par\medskip
\renewcommand{\arraystretch}{1.6}
\begin{tabular}{|l|c|c|}
\hline
\textbf{Gap} & \textbf{On the map (cm)} & \textbf{Real (m)}\\ \hline
First door to third door & & \\ \hline
\end{tabular}
\end{tryit}

\begin{worked}{Practice numbers}
Doors at 1~cm and 5~cm on the map. Gap $5-1=4$~cm, so $4\times2=8$~m.\\
Check on the real hallway: doors at 2~m and 10~m, a gap of $10-2=8$~m.
\end{worked}

\clearpage

\begin{puzzle}{Map to real}
At 1~cm to 2~m, 3~cm on the map is \wline[1.6cm]~m on the ground.
\end{puzzle}

\begin{puzzle}{Real to map}
At 1~cm to 4~m, a real distance of 8~m is \wline[1.6cm]~cm on the map.
\end{puzzle}

\begin{puzzle}{Too small}
A window is 50~cm wide. At 1~cm to 1~m it is \wline[1.6cm]~cm. At 1~cm to 2~m it is \wline[1.6cm]~cm.
\end{puzzle}

\begin{puzzle}{Pick a scale}
You want a poster-sized map of one classroom. Which fits better, 1~cm to 10~cm or 1~cm to 10~m? \wline[3cm]
\end{puzzle}

\begin{puzzle}{Detail to keep}
A fire-drill map must keep some things. Write K for must keep or D for can drop.\\
(a) exit doors \wline[1cm] \quad (b) colour of the carpet \wline[1cm] \quad (c) stair places \wline[1cm] \quad (d) pictures on walls \wline[1cm]
\end{puzzle}

\clearpage

\begin{puzzle}{Scale bar}
A map's scale bar is 2~cm long and stands for 4~m. A copier makes the map 1.5 times bigger. How long is the bar now? \wline[1.6cm]~cm Does it still tell the truth? \wline[2cm]
\end{puzzle}

\begin{puzzle}{Unfair claim}
Ana's map of the hallway at 1~cm to 1~m is 12~cm long. Ben's at 1~cm to 2~m is 6~cm long. Ana says her hallway is longer. What is wrong? \wline[6cm]
\end{puzzle}

\begin{puzzle}{Draw the circles}
A hallway is 12~m long. Its map at 1~cm to 2~m is \wline[1.2cm]~cm long. Fill in the empty circle.
\begin{center}
\begin{tikzpicture}
  \node[bigcir,minimum size=2.4cm] (a) at (0,0) {}; \node[font=\sffamily\scriptsize] at (0,0) {real 12 m};
  \node[cir,dashed] (b) at (4,0) {}; \draw[popA] (a) -- node[above,note] {1 cm to 2 m} (b);
\end{tikzpicture}
\end{center}
\end{puzzle}

\begin{puzzle}{Wrap}
A student wrote ``The map shows everything in the hallway.'' Wrap it with a better claim. \wline[6cm]
\end{puzzle}

\begin{puzzle}{Design a path map}
Pick a path from your door to a place you visit. Say your scale and what you would have to drop. \wline[6cm]
\end{puzzle}
\clearpage

\thispagestyle{fancy}
\subhead{Map record}
Map three paths at the scales shown. Write the real length first.\par\medskip
\renewcommand{\arraystretch}{3.0}
\begin{tabularx}{\linewidth}{|>{\sffamily\small\bfseries}p{2.6cm}|X|X|X|X|}
\hline
Path & Real (m) & Scale & Map (cm) & Checked back? \\ \hline
 & & & & \\ \hline
 & & & & \\ \hline
 & & & & \\ \hline
\end{tabularx}

\vspace{8pt}
Which scale fit best on your page? \wline[7cm]\\[10pt]
What had to be dropped to make it fit? \wline[5cm]\\[10pt]
Sketch one map and add its scale bar.\par\medskip
\begin{tikzpicture}\draw[ink!70,line width=1pt,rounded corners=3pt] (0,0) rectangle (11.4,5.2);\end{tikzpicture}
\clearpage
